\documentclass{resume} % Use the custom resume.cls style

\usepackage[left=0.4in,top=0.4in,right=0.4in,bottom=0.4in]{geometry}
\usepackage[utf8]{inputenc}
\usepackage{cmap}
\usepackage[T2A]{fontenc}
\usepackage[russian]{babel}
\usepackage{graphicx}
\usepackage{xcolor}
\usepackage{tabularx}

% Приглушённый цвет ссылок вместо стандартного насыщенного синего
\definecolor{linkcolor}{RGB}{0, 0, 0}
\hypersetup{colorlinks=true, urlcolor=linkcolor, linkcolor=linkcolor}

\name{Караганов Павел}
\address{Machine Learning}
\address{
    \href{mailto:karaganovpavel119@yandex.ru}{karaganovpavel119@yandex.ru} \hspace{2ex}
    \href{https://t.me/Antichnoste}{@Antichnoste} \hspace{2ex}
    \href{https://github.com/Antichnoste}{github.com/Antichnoste}
}

\begin{document}

\begin{rSection}{О себе}
Студент 3 курса ИТМО. Развиваюсь в сфере машинного обучения и Data Science. Обладаю алгоритмическим фундаментом и интересом к математике. Реализовывал ML-алгоритмы с нуля и обучал нейросети на PyTorch. Хочу применять математические модели и Deep Learning архитектуры для решения реальных продуктовых задач.
\end{rSection}

%----------------------------------------------------------------------------------------
%	EDUCATION SECTION
%----------------------------------------------------------------------------------------
\begin{rSection}{Образование}
{\bf Программная Инженерия}, бакалавриат ИТМО (Факультет ПИиКТ) \hfill {2024 - 2028}
\begin{itemize}
    \itemsep -3pt {} 
    \item Поступил без вступительных испытаний (БВИ) по результатам перечневых олимпиад.
    \item Профильные дисциплины: Алгоритмы и структуры данных, Математический анализ, Линейная алгебра.
\end{itemize}
\end{rSection}

%----------------------------------------------------------------------------------------
% TECHINICAL STRENGTHS	
%----------------------------------------------------------------------------------------
\begin{rSection}{Технологии}
\begin{tabular}{ @{} >{\bfseries}l @{\hspace{4ex}} l }
ML/DL: & scikit-learn, LightGBM, PyTorch, NumPy, pandas \\
Визуализация: & Matplotlib, Seaborn \\
Языки: & Python, C++, SQL \\
Инструменты и ОС: & Jupyter Notebook, Git, Linux, Bash, Docker \\
Иностранный язык: & Английский — B2
\end{tabular}
\end{rSection}

%----------------------------------------------------------------------------------------
%	ACHIEVEMENTS SECTION
%----------------------------------------------------------------------------------------
\begin{rSection}{Олимпиады и Достижения}
\begin{itemize}
    \itemsep -3pt {}
    \item \textbf{Спортивное программирование:} Призёр перечневых олимпиад, участник регионального этапа ВсОШ.
    \item \textbf{Шахматы:} 1-й взрослый разряд (рейтинг РШФ 2000).
\end{itemize}
\end{rSection} 

%----------------------------------------------------------------------------------------
%	WORK EXPERIENCE / PROJECTS SECTION
%----------------------------------------------------------------------------------------
\begin{rSection}{Проекты}

\item \textbf{\href{https://github.com/Antichnoste/ml-from-scratch-to-pytorch}{ML from Scratch to PyTorch}} \hfill \textit{NumPy, PyTorch, scikit-learn}
\begin{itemize}
    \itemsep -3pt {}
    \item Реализовал градиентный спуск, линейную и логистическую регрессию с нуля на NumPy; сравнил со \textbf{sklearn} — расхождение MSE $\sim$0.4\%, ROC-AUC отличается в третьем знаке после запятой.
    \item Спроектировал и обучил \textbf{CNN} на PyTorch для классификации Fashion-MNIST, достиг \textbf{91.99\%} accuracy на test set.
    \item Провёл анализ ошибок модели по confusion matrix: выявил кластер визуально похожих классов (Shirt/\allowbreak T-shirt/\allowbreak Pullover/\allowbreak Coat) как основной источник ошибок.
\end{itemize}

\item \textbf{\href{https://github.com/Antichnoste/store-sales-forecasting}{Store Sales Forecasting}} \hfill \textit{LightGBM, pandas, временные ряды}
\begin{itemize}
    \itemsep -3pt {}
    \item Построил direct multi-horizon пайплайн прогнозирования продаж для \textbf{3\,000\,888 наблюдений}, 54 магазинов, 33 товарных категорий и горизонта 16 дней.
    \item Реализовал \textbf{3-fold walk-forward CV без утечки данных} — лаги строил относительно момента начала каждого 16-дневного прогноза; добавил признаки промоакций, календаря, праздников, магазинов и цен на нефть.
    \item Обучил глобальный и 33 специализированных по товарным категориям \textbf{LightGBM}; OOF-блендинг снизил Kaggle RMSLE с \textbf{1.13406 до 0.42880} (\textminus 62\%).
\end{itemize}

\end{rSection} 

\end{document}
